\bookchapter{第四章：安装和管理App——手机的“技能包”}

\chapterintro{App让手机无所不能，但也要管好它们。就像你会给新员工安排工位、分配权限一样，安装App也需要一套规范流程。}

\booksection{一、官方商店：最安全的“应用超市”}

\begin{itemize}
  \item \textbf{安卓}：Google Play（国外）或各品牌自带商店（如华为应用市场、小米应用商店）。
  \item \textbf{iOS}：只有 App Store。
  \item \textbf{鸿蒙}：华为应用市场。
\end{itemize}

核心安全意义：经过审核、更新有保障、支付安全。行动指南：\textbf{优先只用官方商店下载！}

\booksection{二、纯净/防护模式：开启手机的“防毒墙”}

主流国产品牌（华为、小米、OPPO、vivo、荣耀）都内置了\textbf{纯净模式}或\textbf{防护模式}，这是系统级别的安全开关。\textbf{强烈建议所有用户都开启}。

\booksubsection{它到底是什么？}

纯净模式（也叫“防护模式”）的\textbf{核心功能是阻断第三方 App 的安装}：当它开启时，系统会\textbf{只允许从官方应用商店安装 App}，从浏览器、二维码、聊天软件等渠道下载的安装包会被直接拦截，从源头杜绝恶意软件。

\booksubsection{适用人群}

\begin{itemize}
  \item \textbf{老年用户}：分不清"官方"和"假冒"广告，纯净模式是最稳妥的护身符。
  \item \textbf{对应用权限不敏感的用户}：不需要安装小众 App、测试版、外服游戏。
  \item \textbf{希望提升设备安全性的所有用户}：哪怕是发烧友，开着也比不开强——真有临时需要，再单独关一下即可。
\end{itemize}

\booksubsection{开启路径（以华为/小米为例）}

\begin{itemize}
  \item \textbf{华为/鸿蒙}：[设置] > [系统和更新] >\textbf{[纯净模式]}，选择"开启"。
  \item \textbf{小米}：[设置] > [隐私保护] >\textbf{[保护隐私]}或 [设置] > [应用设置] > [应用安装]，关闭"安装未知应用"。
  \item \textbf{OPPO/vivo}：[设置] > [安全] >\textbf{[安装外部来源应用]}，逐个关闭授权。
\end{itemize}

\booksubsection{应用商店的安全检测机制}

应用商店不是“摆设”，它背后有完整的安全流水线：

\begin{itemize}
  \item \textbf{病毒扫描}：每个 App 上架前会经过多家杀毒引擎联合扫描，识别已知木马和病毒特征。
  \item \textbf{权限审核}：商店会审查 App 申请的权限是否"合理"——一个手电筒 App 要通讯录权限？直接驳回或下架。
  \item \textbf{开发者认证}：上架需要企业资质、个人身份证或信用卡认证。出问题能找到人，不是"黑 App"。
  \item \textbf{持续监控}：已上架的 App 也会被持续复扫，新发现的恶意行为会被强制下架。
\end{itemize}

\booksection{三、第三方来源与“电子寄生虫”：风险极高！}

从网页或不明链接下载App，可能夹带病毒、窃取隐私、盗刷支付。这里我们特别提一个新概念：\textbf{“电子寄生虫”}。

\textbf{什么是“电子寄生虫”？}\par

指的是那些\textbf{被植入木马的“破解版”软件}——你以为是“省了钱”，实际是“中了招”。常见形式：

\begin{itemize}
  \item "破解版"、"免费版"、"去广告版"付费软件；
  \item "刷赞"、"抢红包"、"游戏外挂"等灰色工具；
  \item 群聊里"小广告"推荐的所谓"内部版"App。
\end{itemize}

这些程序表面能用，背后可能在\textbf{偷通讯录、读短信、监听麦克风、后台挖矿}。纯净模式是防范它们\textbf{第一道防线}——直接不让你装。

\booksubsection{第二道防线建议：}

\begin{itemize}
  \item \textbf{定期全盘病毒扫描}：使用手机自带安全中心（如华为手机管家、小米安全中心）每月扫一次。
  \item \textbf{谨慎授予敏感权限}：通讯录、短信、位置、相机、麦克风——这五项"高危权限"必须严格审批。
  \item \textbf{关注耗电与流量异常}：手机突然发烫、电量掉得快、流量半夜跑——八成是后台藏了"寄生虫"。
\end{itemize}

如果万不得已要装第三方 App：

\begin{enumerate}
  \item 确保来源绝对可信。
  \item 安装前仔细看权限要求。
  \item 事后用安全App扫描。
\end{enumerate}

终极建议：在系统设置里，关掉“安装未知应用”的权限。

\booksection{四、装、更新、卸载}

\begin{itemize}
  \item \textbf{安装}：商店里点"获取"或"安装"，装完首次打开会问你要权限。
  \item \textbf{更新}：\textbf{必须更！}主要是修复安全漏洞。建议开启自动更新。
  \item \textbf{卸载}：长按图标拖到卸载区。想清得更彻底，去 [设置] > [应用管理] 里找到并卸载。
\end{itemize}

\booksection{五、权限管理：你是App的“老板”}

App要权限（比如通讯录、位置）有时合理（地图要位置），有时不合理（手电筒要通讯录）。

\begin{itemize}
  \item \textbf{权限选项}："始终允许"、"仅在使用时允许"、"禁止"、"每次询问"。
  \item \textbf{管理原则}：
  \begin{enumerate}
    \item 最小化：非必要不给。
    \item 事后审计：定期去权限设置里看看谁在"越权"。
    \item 情境化：像打车软件，只在用时给位置权限。
  \end{enumerate}
\end{itemize}

\begin{tcbwarn}{重点盯防的权限}
\begin{itemize}
  \item 通讯录/短信：可能被用来发垃圾短信。
  \item 实时位置：可能被追踪行踪。
  \item 相机/麦克风：可能被偷拍偷录。
  \item 文件访问：可能读取你的隐私文件。
\end{itemize}

\end{tcbwarn}

\booksection{六、软件安装包的构成要素}

你有没有想过，从应用商店下载的那个安装包，里面到底装了什么？了解安装包的构成，能帮你更好地理解App是怎么工作的，也能在遇到问题时知道该从哪里排查。

\booksubsection{安卓APK文件结构}

APK（Android Package Kit）本质是一个ZIP压缩包，解压后可以看到以下核心目录和文件：

\begin{itemize}
  \item \textbf{AndroidManifest.xml}：应用的"身份证"和"说明书"。记录了应用包名、版本号、所需权限、支持的系统版本、入口Activity等核心信息。系统安装应用时首先读取这个文件。
  \item \textbf{classes.dex}：应用的核心代码。Java/Kotlin源代码编译后变成Dalvik字节码，存在这个文件里。一个应用可能有多个dex文件（如classes2.dex、classes3.dex）。
  \item \textbf{resources.arsc}：编译后的资源索引表。记录了所有资源的ID和对应关系。
  \item \textbf{res/目录}：资源文件夹。包含图片、布局文件、字符串、颜色定义等。不同分辨率的图片放在不同的drawable文件夹里（如drawable-hdpi、drawable-xhdpi）。
  \item \textbf{assets/目录}：原始资源文件夹。存放不经过编译的资源，如网页文件、字体、音频等。App运行时可以通过AssetManager读取。
  \item \textbf{lib/目录}：本地库文件夹。存放编译好的.so文件（C/C++代码编译产物），按CPU架构分类（如arm64-v8a、armeabi-v7a）。游戏和需要高性能计算的App通常有这个目录。
  \item \textbf{META-INF/目录}：签名信息文件夹。包含MANIFEST.MF、CERT.SF、CERT.RSA三个文件，用于验证APK的完整性和签名。
\end{itemize}

\booksubsection{iOS IPA文件结构}

IPA（iOS App Store Package）同样是压缩包，结构略有不同：

\begin{itemize}
  \item \textbf{Payload/目录}：核心内容所在。里面有一个.app文件夹，包含应用的所有内容。
  \item \textbf{Info.plist}：应用的配置文件，类似安卓的AndroidManifest.xml。记录了应用ID、版本、权限需求、支持的设备等。
  \item \textbf{可执行文件}：与.app文件夹同名的主程序，是编译后的Mach-O格式二进制文件。
  \item \textbf{Frameworks/目录}：应用依赖的动态库和框架。
  \item \textbf{资源文件}：图片、音频、视频等资源，通常存放在.app文件夹内。
  \item \textbf{\_CodeSignature/目录}：代码签名信息，用于验证应用的完整性和来源。
\end{itemize}

\booksubsection{签名验证机制}

为什么从官方商店下载的App更安全？签名验证是关键。

\begin{itemize}
  \item \textbf{安卓签名}：开发者用私钥对APK进行签名，系统安装时用公钥验证。如果APK被篡改（比如植入病毒），签名就会失效，系统会拒绝安装。Google Play还要求App Signing，由Google保管签名密钥，进一步保护应用安全。
  \item \textbf{iOS签名}：苹果的审核更严格。每个IPA必须有苹果颁发的证书签名，否则无法安装。这就是为什么iOS只能从App Store或企业证书安装应用。企业证书签名的应用有效期有限，且苹果可以随时吊销。
  \item \textbf{签名的作用}：
  \begin{enumerate}
    \item 验证来源：确认应用确实来自声称的开发者。
    \item 验证完整性：确保应用在传输过程中未被篡改。
    \item 防止冒充：恶意应用无法伪装成其他应用。
  \end{enumerate}
\end{itemize}

\booksubsection{安装包大小与优化}

你可能会发现，应用商店显示的App大小和安装后占用的空间不一样。这是因为：

\begin{itemize}
  \item \textbf{下载大小}：压缩后的APK/IPA大小。
  \item \textbf{安装大小}：解压后的大小，通常比下载大小大很多。
  \item \textbf{运行时占用}：App运行时还会产生缓存、数据文件，占用更多空间。
\end{itemize}

现代App的优化手段：

\begin{itemize}
  \item \textbf{App Bundle（安卓）}：Google Play的动态分发技术。用户下载时只获取适合自己设备的资源（如只下载对应屏幕分辨率的图片），大大减少下载大小。
  \item \textbf{App Thinning（iOS）}：类似的技术，苹果会根据设备自动裁剪不需要的资源。
  \item \textbf{按需下载}：游戏等大型应用可以先下载核心包，后续再下载资源包。
\end{itemize}

\begin{tcbinfo}{小知识：为什么有的App安装包特别大？}
游戏类App通常最大，因为包含大量高清贴图、3D模型、音频文件。社交类App也比较大，因为要支持各种功能模块。工具类App通常最小。如果你发现某个工具类App特别大，要警惕是否捆绑了不必要的功能。

\end{tcbinfo}

\begin{tcbpractice}{实操环节：审查3个App的权限}
选3个常用App（如社交、购物、工具类），进入 [设置] 的应用权限管理页，逐项审查，把不必要的敏感权限关掉或改成“仅在使用时允许”。

\end{tcbpractice}
